\section{Round-thirty-one reconstruction: buffered current response and covariance on a common Hilbert level}
\label{sec:r31-a1}

The current derivative genuinely loses regularity.  The covariance theorem is
therefore formulated on one fixed reporting space containing every derivative
that enters it.  This removes the type mismatch identified in the
Round-Thirty report without suppressing the parameter response.  We also
verify, for the affine seam observable itself, the two-sided projection
summability required by the functional limit theorem.

\subsection{Loss scale and a common reporting level}

Let $\widehat\Phi^s$ be the closed compatible test space on the labelled
resolution and put $\mathbb J^s=(\widehat\Phi^s)'$.  For $s'>s$ there is a
continuous dense inclusion
\[
 \iota_{s,s'}:\mathbb J^s\hookrightarrow\mathbb J^{s'}.
\]
Material differentiation consumes two test derivatives, hence
\[
 \nabla_a:\mathbb J^s\longrightarrow\mathbb J^{s+2}.      \tag{A1.1}
\]
Fix an integer $r\ge1$ and define the common reporting Hilbert space
\[
 \mathscr H_r:=\mathbb J^{m+2r}.                           \tag{A1.2}
\]
A current jet $J_a^{(j)}\in\mathbb J^{m+2j}$ is always viewed in
$\mathscr H_r$ through $\iota_{m+2j,m+2r}$.

\begin{lemma}[Buffered differentiation]
\label{lem:r31-a1-buffer}
Suppose $a\mapsto J_a$ has a loss-indexed jet through order $r$ and
\[
 \sup_a\sum_{j=0}^r
 \norm{\iota_{m+2j,m+2r}J_a^{(j)}}_{\mathscr H_r}<\infty. \tag{A1.3}
\]
Then $a\mapsto \iota_{m,m+2r}J_a$ is $C^r$ as an
$\mathscr H_r$-valued map and
\[
 {\dd^j\over\dd a^j}\iota_{m,m+2r}J_a
 =\iota_{m+2j,m+2r}J_a^{(j)}.                              \tag{A1.4}
\]
No derivative is asserted in the smaller space $\mathbb J^m$.
\end{lemma}

\begin{proof}
For a transported test $\phi\in\widehat\Phi^{m+2r}$, finite-cylinder
material differentiation and transposition give (A1.4) weakly.  The uniform
bound (A1.3) makes the difference quotients weakly bounded in the Hilbert
space.  The projective tail estimate below makes them Cauchy in norm.  The
Hilbert-space fundamental theorem of calculus then upgrades weak derivatives
to strong derivatives.  The moving Dirac model is consistent with the
statement: $\partial_a\delta_a$ need not belong to $\mathbb J^m$, but both
$\delta_a$ and all derivatives through order $r$ belong to the common larger
space $\mathscr H_r$.
\end{proof}

\subsection{Coherent shell derivatives}

Let $\mu_a=\bigotimes_{k\in\Z}w(a)$ on
$\Sigma=\{1,2,3,4\}^{\Z}$ and let $\cF_N$ be generated by
$\omega_k$, $|k|\le N$.  A coherent current jet consists of martingales
$J^{(j)}_{a,N}$ and boundary shells $b^{(j)}_{a,q}$ satisfying
\[
 E_a[J^{(j)}_{a,N+1}\mid\cF_N]=J^{(j)}_{a,N},             \tag{A1.5}
\]
\[
 \partial_aJ^{(j)}_{a,N}=J^{(j+1)}_{a,N}
   +\sum_{q\le N}b^{(j)}_{a,q}.                            \tag{A1.6}
\]
For some $\eta>0$ require
\begin{equation}\tag{A1.7}
 \sum_{j=0}^r\sup_a\left[
  \sup_N\norm{J^{(j)}_{a,N}}_{L^2(\mathscr H_r)}
 +\sum_{N\ge0}e^{\eta N}
  \norm{J^{(j)}_{a,N+1}-J^{(j)}_{a,N}}_{L^2(\mathscr H_r)}
 +\sum_{q\ge0}e^{\eta q}
  \norm{b^{(j)}_{a,q}}_{L^2(\mathscr H_r)}
 \right]<\infty.
\end{equation}
The cumulative boundary current is
$B_a^{(j)}=\sum_qb^{(j)}_{a,q}$; it is the shell, not the cumulative sum,
that decays.

\begin{theorem}[Loss-buffered projective response]
\label{thm:r31-a1-response}
Under (A1.7), $J^{(j)}_{a,N}$ converges in
$L^2(\mu_a;\mathscr H_r)$ to $J_a^{(j)}$.  The map
$a\mapsto E_aJ_a^{(0)}$ is $C^r$ in $\mathscr H_r$.  Its first derivative is
\begin{equation}\tag{A1.8}
 {\dd\over\dd a}E_aJ_a^{(0)}
 =E_a[J_a^{(1)}+B_a^{(0)}]
 +\sum_{k\in\Z}E_a\left[
   \bigl(J_a^{(0)}-E_a[J_a^{(0)}\mid\cF_{|k|-1}]\bigr)
   \partial_a\log w_{\omega_k}(a)\right],
\end{equation}
with absolute convergence in $\mathscr H_r$.  Higher derivatives are the
normally convergent Bell expansions generated by the same shells.
\end{theorem}

\begin{proof}
The first series in (A1.7) is a summable martingale-difference bound and the
second is a normally convergent boundary-shell bound.  On a score shell of
width $q$, the product score has $L^2$ norm $O(q^{1/2})$; multiplying by the
$O(e^{-\eta q})$ current shell gives a summable majorant.  Repeated
parameter differentiation produces only a fixed polynomial in $q$, still
absorbed by $e^{-\eta q}$.  Finite-product differentiation, followed by the
projective limit, proves (A1.8).  Lemma~\ref{lem:r31-a1-buffer} fixes the
common target space at every order.
\end{proof}

\subsection{The affine seam current satisfies the projection criterion}

For a depth-$n$ symbolic word $\omega$, let $p_\omega(a)$ be its branch
weight.  The affine parent-fold geometry gives, for $j\le r$,
\[
 \norm{\iota_{m+2j,m+2r}
   \nabla_a^j[J_a^{\rm seam}]_\omega}_{\mathscr H_r}
 \le C_j(1+n)^jp_\omega(a),                               \tag{A1.9}
\]
and every incidence row has multiplicity at most $d_0$.  Assume
\[
 \vartheta^2:=\sup_a\sum_{i=1}^4w_i(a)^2<1,
 \qquad \beta d_0^2\vartheta^2<1.                         \tag{A1.10}
\]

\begin{proposition}[Seam shell and two-sided projection decay]
\label{prop:r31-a1-seam}
The affine seam current satisfies (A1.7).  If
$Y_k=J_a^{(0)}\circ\sigma^k-E_aJ_a^{(0)}$ and
\[
 P_0Y_q=E[Y_q\mid\cF_0^-]-E[Y_q\mid\cF_{-1}^-],
\]
then, for some $\theta<1$,
\[
 \norm{P_0Y_q}_{L^2(\mathscr H_r)}
 \le C(1+|q|)^r\theta^{|q|},                              \tag{A1.11}
\]
so in particular
\[
 \sum_{q\in\Z}(1+|q|)^{1/2}
 \norm{P_0Y_q}_{L^2(\mathscr H_r)}<\infty.                \tag{A1.12}
\]
\end{proposition}

\begin{proof}
Squaring (A1.9), summing all words of length $n$, and using the product
identity gives
\[
 \sum_{|\omega|=n}\norm{\nabla_a^j[J^{\rm seam}]_\omega}^2
 \le C_j(1+n)^{2j}\vartheta^{2n}.
\]
The incidence weight is absorbed by (A1.10).  The difference between depths
$n-1$ and $n$ is exactly the new seam shell, proving (A1.7).  The projection
$P_0Y_q$ can only see seam sheets whose symbolic ancestry crosses both the
origin and coordinate $q$; their minimum depth is $|q|$.  Summing the same
product estimate from that depth onward yields (A1.11).  This proves the
projection condition for the motivating observable instead of imposing it as
an unrelated hypothesis.
\end{proof}

\subsection{Functional limit and covariance response}

Let $D_0=\sum_{q\in\Z}P_0Y_q$ and $D_k=D_0\circ\sigma^k$.  Define the genuine
polygonal interpolation
\[
 \widetilde W_n(t)=n^{-1/2}\left[
   \sum_{k<\lfloor nt\rfloor}Y_k
  +(nt-\lfloor nt\rfloor)Y_{\lfloor nt\rfloor}\right].   \tag{A1.13}
\]

\begin{theorem}[Common-level trace-class FCLT and covariance response]
\label{thm:r31-a1-fclt}
The processes $\widetilde W_n$ converge in
$C([0,1];\mathscr H_r)$ to Brownian motion with covariance
\[
 Q_a=E_a[D_0\otimes D_0]\in\cS_1(\mathscr H_r).           \tag{A1.14}
\]
Moreover $a\mapsto Q_a$ is $C^{r-1}$ in trace norm on the single space
$\mathscr H_r$, and
\[
 \partial_aQ_a
 =E_a[D_0'\otimes D_0+D_0\otimes D_0']
  +E_a[(D_0\otimes D_0)S_a],                              \tag{A1.15}
\]
where the last term is the normally convergent centred projective score
series.  Every vector in (A1.15), including $D_0'$, belongs to
$\mathscr H_r$.
\end{theorem}

\begin{proof}
Equation (A1.12) gives the Hilbert Maxwell--Woodroofe maximal approximation;
the remainder is $o_{L^2}(\sqrt n)$ uniformly in the maximum over partial
sums.  The martingale FCLT for $(D_k)$ and the maximal-jump estimate transfer
the limit to (A1.13).  Parseval gives
$\operatorname{tr}Q_a=E_a\norm{D_0}_{\mathscr H_r}^2$.

Differentiate the projection series using
Theorem~\ref{thm:r31-a1-response}.  Proposition~\ref{prop:r31-a1-seam}
gives an exponentially summable bound for every derivative through order
$r-1$.  Since
$\norm{x\otimes y}_{\cS_1}=\norm x\norm y$, the differentiated rank-one
series is normally convergent in trace norm.  Crucially, the covariance is
formed on $\mathscr H_r$, the common level containing all derivatives; no
trace-class derivative is claimed on the smaller space $\mathbb J^m$.
\end{proof}

This closes both Round-Thirty objections: the covariance differentiates on a
correctly buffered Hilbert level, and the seam current itself satisfies the
weighted two-sided projection estimate used by the FCLT.
